% Header aus der Vorlage
\documentclass[a4paper,11pt, footheight=26pt]{scrartcl}
\usepackage[head=23pt]{geometry}	% head=23pt umgeht Fehlerwarnung, dafür größeres "top" in geometry
\geometry{a4paper, top=30mm, left=20mm, right=20mm, bottom=22mm,headsep=10mm, footskip=12mm}

% Input inkl. Umlaute, Silbentrennung
\usepackage[T1]{fontenc}
\usepackage[utf8]{inputenc}
\usepackage[ngerman]{babel}
\usepackage{csquotes}	% Anführungszeichen

% HTW Corporate Design: Arial (Helvetica)
\usepackage{helvet}
\renewcommand{\familydefault}{\sfdefault}

% Style-Aufhübschung
\usepackage{soul, color}	% Kapitälchen, Unterstrichen, Durchgestrichen usw. im Text
\usepackage{scrlayer-scrpage}	% Kopf-/Fußzeile
%\usepackage{titleref}
\usepackage[perpage]{footmisc}	% Fußnotenzählung Seitenweit, nicht Dokumentenweit

% Mathe usw.
\usepackage{amssymb}
\usepackage[fleqn]{amsmath}	% fleqn: align-Umgebung rechtsbündig
\usepackage{xcolor}
\usepackage{esint}	% Schönere Integrale, \oiint vorhanden
\everymath=\expandafter{\the\everymath\displaystyle}	% Mathe Inhalte werden weniger verkleinert
\usepackage{wasysym}	% mehr Symbole, bspw \lightning

% tikz usw.
\usepackage{tikz}
\usepackage{pgfplots}
\pgfplotsset{compat=1.11}	% Umgeht Fehlermeldung
\usetikzlibrary{graphs}
%\usetikzlibrary{through}	% ???
\usetikzlibrary{arrows}
\usetikzlibrary{arrows.meta}	% Pfeile verändern / vergrößern: \draw[-{>[scale=1.5]}] (-3,5) -> (-3,3);
\usetikzlibrary{automata,positioning} % Zeilenumbruch im Node node[align=center] {Text\\nächste Zeile}
\usetikzlibrary{patterns}	% Schraffierte Füllung
\tikzstyle{reverseclip}=[insert path={	% Inverser Clip \clip
	(current page.north east) --
	(current page.south east) --
	(current page.south west) --
	(current page.north west) --
	(current page.north east)}
% Nutzen: 
%\begin{tikzpicture}[remember picture]
%\begin{scope}
%\begin{pgfinterruptboundingbox}
%\draw [clip] DIE FLÄCHE, IN DER OBJEKT NICHT ERSCHEINEN SOLL [reverseclip];
%\end{pgfinterruptboundingbox}
%\draw DAS OBJEKT;
%\end{scope}
%\end{tikzpicture}
]	% Achtung: dafür muss doppelt kompliert werden!
\usepackage{graphpap}	% Grid für Graphen

% Tabular
\usepackage{longtable}	% Große Tabellen über mehrere Seiten
\usepackage{multirow}	% Multirow/-column: \multirow{2[Anzahl der Zeilen]}{*[Format]}{Test[Inhalt]} oder \multicolumn{7[Anzahl der Reihen]}{|c|[Format]}{Test2[Inhalt]}
\renewcommand{\arraystretch}{1.3} % Tabellenlinien nicht zu dicht
\usepackage{colortbl}
\arrayrulecolor{gray}	% heller Tabellenlinien

% Nützliches
\usepackage{verbatim}	% u.a. zum auskommentieren via \begin{comment} \end{comment}
\usepackage{tabto}	% Tabs: /tab zum nächsten Tab oder /tabto{.5 \CurrentLineWidth} zur Stelle in der Linie
\NumTabs{6}	% Anzahl von Tabs pro Zeile zum springen
\usepackage{listings} % Source-Code mit Tabs
\usepackage{lstautogobble} 
\usepackage{enumitem}	% Anpassung der enumerates
\setlist[enumerate,1]{label=\arabic*.)}	% global andere Enum-Items
\usepackage{letltxmacro} % neue Definiton von Grundbefehlen
% Nutzen:
%\LetLtxMacro{\oldemph}{\emph}
%\renewcommand{\emph}[1]{\oldemph{#1}}

% Einrichtung von lst
\lstset{
basicstyle=\ttfamily, 
mathescape=true, 
escapeinside=||, 
autogobble, 
tabsize=2,
basicstyle=\footnotesize\sffamily\color{black},
frame=single,
rulecolor=\color{lightgray},
numbers=left,
numbersep=5pt,
%basicstyle=\footnotesize\sffamily\color{black},
commentstyle=\color{gray},
%frame=single,
%numbers=left,
%numbersep=5pt,
numberstyle=\tiny\color{gray},
keywordstyle=\color{green},
%showspaces=false,
showstringspaces=false,
stringstyle=\color{orange},
tabsize=2,
literate=%
    {Ö}{{\"O}}1
    {Ä}{{\"A}}1
    {Ü}{{\"U}}1
    {ß}{{\ss}}1
    {ü}{{\"u}}1
    {ä}{{\"a}}1
    {ö}{{\"o}}1
    {~}{{\textasciitilde}}1
}

% BibTeX
\usepackage[backend=bibtex, bibencoding=ascii]{biblatex}	% BibTeX
\usepackage{makeidx}
%\makeglossary
%\makeindex

% Grafiken
\usepackage{graphicx}
\usepackage{epstopdf}	% eps-Vektorgrafiken einfügen

% pdf-Setup
\usepackage{pdfpages}
\usepackage[bookmarks,%
bookmarksopen=false,% Klappt die Bookmarks in Acrobat aus
colorlinks=true,%
linkcolor=black,%
citecolor=red,%
urlcolor=green,%
]{hyperref}

% Titel, Autor usw. werden vor dem Anfang des Dokuments in einem Rutsch definiert…
\newcommand{\DTitel}[1]{\newcommand{\Dokumententitel}{#1}}
\newcommand{\DUntertitel}[1]{\newcommand{\Dokumentenuntertitel}{#1}}
\newcommand{\DAutor}[1]{\newcommand{\Dokumentenautor}{#1}}
\newcommand{\DNotiz}[1]{\newcommand{\Dokumentennotiz}{#1}}
% … Deswegen folgendes erst Nach Dokumentenbeginn ausführen:
\AtBeginDocument{
	\hypersetup{
		pdfauthor={\Dokumentenautor},
		pdftitle={HTW Dresden | \Dokumententitel - \Dokumentenuntertitel},
	}
	\automark[section]{section}
	\automark*[subsection]{subsection}
	\pagestyle{scrheadings}
	\ihead{\includegraphics[height=1.7em]{../LaTeX_master/HTW-Logo.eps}}
	\ohead{\Dokumententitel}
	\ofoot{\footnotesize{\textcolor{darkgray}{Mitschrift von\\ \Dokumentenautor}}}
	% Titelseite
	\title{\includegraphics[width=0.35\textwidth]{../LaTeX_master/HTW-Logo.eps}\\\vspace{0.5em}
	\Huge\textbf{\Dokumententitel} \\\vspace*{0,5cm}
	\Large \Dokumentenuntertitel \\\vspace*{4cm}}
	\author{\textcolor{darkgray}{Mitschrift von \Dokumentenautor} \vspace*{1cm}\\\Dokumentennotiz}
}

%% EIGENE BEFEHLE

%Farbdefinitionen
\definecolor{red}{RGB}{180,0,0}
\definecolor{green}{RGB}{75,160,0}
\definecolor{blue}{RGB}{0,75,200}
\definecolor{orange}{RGB}{255,128,0}
\definecolor{yellow}{RGB}{255,245,0}
\definecolor{purple}{RGB}{75,0,160}
\definecolor{cyan}{RGB}{0,160,160}
\definecolor{brown}{RGB}{120,60,10}

% Textfarbe ändern
\newcommand{\tred}[1]{\textcolor{red}{#1}}
\newcommand{\tgreen}[1]{\textcolor{green}{#1}}
\newcommand{\tblue}[1]{\textcolor{blue}{#1}}
\newcommand{\torange}[1]{\textcolor{orange}{#1}}
\newcommand{\tyellow}[1]{\textcolor{yellow}{#1}}
\newcommand{\tpurple}[1]{\textcolor{purple}{#1}}
\newcommand{\tcyan}[1]{\textcolor{cyan}{#1}}
\newcommand{\tbrown}[1]{\textcolor{brown}{#1}}

% Umstellen der Tabellen Definition
\newcommand{\mpb}[1][.3]{\begin{minipage}{#1\textwidth}\vspace*{3pt}}
\newcommand{\mpe}{\vspace*{3pt}\end{minipage}}

\newcommand{\resultul}[1]{\underline{\underline{#1}}}
\newcommand{\parskp}{$ $\\}	% new line after paragraph
\newcommand{\corr}{\;\widehat{=}\;}
\newcommand{\mdeg}{^{\circ}}

\newcommand{\nok}[2]{\begin{pmatrix}#1\\#2\end{pmatrix}}	% n über k

\bibliography{../Literatur/HTW_Literatur.bib}

% Definition von Titel, Autor usw.
\DTitel{Programmierung 1}
\DUntertitel{Übungsskript}
\DAutor{Falk-Jonatan Strube}
\DNotiz{Übung von Prof. Dr.-Ing. Beck}

\begin{document}

\maketitle
\newpage
\tableofcontents
\newpage

\section{Eingebaute Datentypen}

\begin{itemize}
\item int
\item string
\item short (oft 16 Bit / 2 Byte)
\item long (so groß wie eine Adresse, abhängig vom System [32/64 bit])
\item char (kann Symbol sein, aber auch eine Zahl [die Symbol represäntiert], 8 Bit)
\item float
\begin{itemize}
\item kann einige Zahlen (ganzzahlig) nicht darstellen (bspw. 2)
\item hat Probleme sehr große und sehr kleine Zahlen miteinander zu addieren (durch Normierung der Exponenten kann die kleine zu 0 werden, oder Nachkommastellen verloren gehen)
\end{itemize}
\item Achtung: boolean kein Datentyp in C!\\
$\curvearrowright$ Abfrage von true/false durch int: \\
$int=0\corr false$\\
$int \not =0 \corr true$
\end{itemize}

\subsection{Zahlentypen}
Zahl: $\underset{8\cdot 10^0}{1} \; \underset{0 \cdot 10^1}{0} \;\underset{1\cdot 10^2}{8}$\\
$\Rightarrow$ 10er-System (Decimal)\\
Zahl: $0110\underset{\leftarrow usw.}{1}\;\underset{1\cdot 2^2}{0}\;\underset{0\cdot 2^1}{0}\;\underset{0\cdot 2^0}{0} = 0 + 0 + 4+8+0+32+64+0 =108$\\
$\Rightarrow$ 2er-System (Binär)\\
Zahl: $\underset{1}{001}|\underset{5}{101}|\underset{4}{100}=108$\\
$\Rightarrow$ 8er-System (Octal)\\
Zahl: $\underset{6}{0110}|\underset{C}{1100}=108$\\
$\Rightarrow$ 16er-System (Hexa)\\
Unterschied: $108_{/10}$, $01101100_{/2}$, $154_{/8}$ (in C gekennzeichnet durch $0154\rightarrow$ Octalzahl) und $6C_{/16}$ (in C gekennzeichnet durch $0x6C$)

\paragraph{Veranschaulichung}
\begin{align*}
108:2&=54 &R \emptyset\\
65:2 &= 27 &R \emptyset\\
27:2&= 13 &R 1\\
13:2 &= 6 &R 1\\
6:2 &=3 &R \emptyset\\
3:2 &= 1 &R 1\\
1:2&= 0 &R 1\\
&\Rightarrow 1101100 \text{ von unten nach oben gelesen}&
\end{align*}

\begin{align*}
108:8 &= 13 &R4\\
13:8 &= 1 &R5\\
1:8 &= 0 &R1\\
&\Rightarrow 154&
\end{align*}

\begin{align*}
108: 16 &= 6 &R12 = RC\\
6:16 &= 0 &R6\\
&\Rightarrow 6C&
\end{align*}

Beispielzahl 0x12AB\medskip\\
Speicherblock:\\
\begin{tabular}{|l|l|l|l|}
\hline
1&2&A&B\\
\hline
\end{tabular} big-endian\\
\begin{tabular}{|l|l|l|l|}
\hline
B&A&1&2\\
\hline
\end{tabular}\\
\begin{tabular}{|l|l|l|l|}
\hline
A&B&1&2\\
\hline
\end{tabular} little-endian\\
Letzte Version ist die, die heutzutage meistens (Intel) verwendet wird: Das niederwertigste Byte liegt auf der niedrigsten Adresse.

\paragraph{2er Komplement}
positive Zahl: $\boxed{0}110 \: 1100$\\
Negation: $1001\: 0011$\\
$+1$\\
Komplement: $1001 \: 0100 = -108 = 0x94$\\
\begin{tabular}{r | l l l l l l l l l}
 & $0$&$1$&$1$&$0$&$0$&$1$&$1$&$0$ & +108\\
$_1$&$1_1$&$0_1$&$0_1$&$1_1$&$0_1$&$1$&$0$&$0$ & -108\\
\hline
$1$&$0$&$0$&$0$&$0$&$0$&$0$&$0$&$0$&\\

\end{tabular}

\subsection{Variable}
4 Kennzeichen einer Variable:
\begin{itemize}
\item Adresse im Speicher (Ort)
\item Datentyp (Verarbeitungsbreite)
\item Bitkombination (Wert)
\item Symbolischer Name
\end{itemize}
Ein Vektor fasst mehrere Variablen gleichen Datentyps unter einer zusammen.

\subsubsection{Initialisierung}
Bei der Initialisierung hat die Variable einen Ausgangswert:
\begin{itemize}
\item Initialisierung innerhalb einer Funktion: zufälliger Wert (alte Speicherbelegung)
\item Init. außerhalb einer Funktion: 0
\end{itemize}

\begin{lstlisting}[language=C]
char c='c';	// 'c'=99 (ASCII)
char c=99;
\end{lstlisting}
Die Anführungszeichen bei der Wertzuweisung einer char ist nur bei einem Zeichen wichtig, nicht bei Zahlen.
\begin{itemize}
\item $'x'$ einzelne Anführungszeichen für ein einzelnes Zeichen
\item $"xy"$ doppelte Anführungszeichen bei mehrere Zeichen
\end{itemize}

\begin{lstlisting}[language=C]
int i=8, j=5, k;
char c=99, d='d';
float x=0.005, y=-0.01, z;
x=(i+j);	// k bleibt int (Wert: 13)
z=y+x;	// x+y werden als double zusammengerechnet (Rechnung immer in double).
				// z ist dann wieder float (Wert: -0.005)
k=x+y;	// k ist int, beim Rechnen wird gebrochenzahliger Teil abgschnitten (Wert: 0)
z=i/j;	// i/j wird in int gerechnet, also 8/5=1 und nicht 1.6! 
				// z ist dann trotzdem float (Wert: 0.000...)
char a,b;
a=b=c;	// erst wird b=c='c'=99, dann wird a=(b=c)=99 zugewiesen. 
				// a und b haben nachfolgend den Wert 'c'
i=j=1.1;	// j=1.1(bzw. 1), also i auch Wert 1
z=k=x;	// k=x=0 (wird abgeschnitten), also z=0.00000...
i+=2;	// i=i+2;
y-=x;	// y=y-x=-0.015
k=j=5?i:j;	// das selbe wie k=((j==5)?i:j);
						// ist j=5? Wenn ja, dann k=i. Wenn nein, dann k=j.
z=y>=0?1.0:2.0;	// das selbe wie z=((y>=0)?1.0:2.0);
printf("%d\n", (3*i-2*j)%(2*d-c) );	// Punkt vor Strich! 14%101=14
printf("%d\n", (2*((i/5)*(4(j-3))%(i+j-2) );
// Achtung: i/5=8/5=1 (s.o.)!
// 2*9=18
printf("%d\n", (i-3*j)%(c+2*d)/(x-y) );
// -7%29 = -7
// dann: -7/0.015)=-466.6667
printf("%d\n", i<=j );	// 0
printf("%d\n", i>=j );	// 1
printf("%d\n", i==6 );	// 0
printf("%d\n", i=!6 );	// !6 entspricht !(ungleich Null)=0
printf("%d\n", i!=6 );	// 1
printf("%d\n", i>0 && j<5 );
printf("%d\n", i&j );	// i bitweise mit j Verknüpft (ge-UND-et):
//  00001000
// &00000101
// =00000000
i=8, j=5;
printf("%s\n",i && j? "true":"false");
printf("%d\n",i && j);
printf("%d\n",i & j);
\end{lstlisting}

\section{Ausdrücke}

Simple Sort

\begin{lstlisting}[language=C]
#include <stdio.h>

int data[] = {7,3,9,2,5};

int main(){
	int ige, iro; // entsprechende Pfeile unter den Zahlen auf Papier
	for (irt=0; irt<(5-1); irt++={
		for (ige = irt+1, ige<5, ige++){
			if (data[ige] < data[irt]){
				int tmp = data [ige];
				data[ige] = data[irt];
				
				data[irt] = tmp;
				// tauschen alternativ (ohne Zwischenspeichern): (^= ist XOR)
				// data[irt]^=data[ige];
				// data[ige]^=data[irt];
				// data[irt]^=data[ige];
			}
		}	
	}
	for (irt=0; irt<5; irt++={
		printf("%d ", data[irt]);
	}
	printf("\n");
	return 0;
}
\end{lstlisting}

Alphabetische Sortierung

\begin{lstlisting}[language=C]
#include <stdio.h>

#define N 10

//[10]: länge der Zeichenkette
char data[][10] = {"Max", "Huckebein", "Bolte", "Lempel", "Maecke", 
	"Helene", "Antonius", "Schlich", "Moritz", "Boeck"};

int main(){
	int ige, iro, ibl; // entspr. Pfeile unter den Zahlen auf Papier
	for (irt=0; irt<(N-1); irt++={
		for (ige = irt+1, ige<N, ige++){
			for (ibl = 0; data[irt][ibl] == data[ige][ibl] && 
				data[irt][ibl]!=0; ibl++){
				;
			}
			if(data[irt][ibl] > data[ige][ibl]){
				char tmp;
				// ibl muss nicht auf 0 gesetzt werden, vertauscht muss sowieso
				// erst ab dem ungleichen Zeichen getauscht werden
				for (/*ibl = 0*/; ibl<N ; ibl++){
					tmp = data[irt][ibl];
					data[irt][ibl] = data[ige][ibl];
					data[ige][ibl] = tmp;
					// alternativ wieder:
					// data[irt][ibl] ^= data[ige][ibl];
					// data[ige][ibl] ^= data[irt][ibl];
					// data[irt][ibl] ^= data[ige][ibl];
			}
		}	
	}
	for (irt=0; irt<N; irt++={
		printf("%d ", data[irt]);
	}
	printf("\n");
	return 0;
}
\end{lstlisting}

\section{Speicherklassen}
\begin{itemize}
\item Register: Prozessor-Register relativ schnell
\item Volatile: Variable wird immer im Hauptspeicher gespeichert (Gegenteil von Register)
\item Static: Liegt die Variable in einer Funktion, dann existiert sie über die gesamte Laufzeit des Programms (kann aber trotzdem nur innerhalb der Funktion verwendet werden). Liegt die Variable außerhalb einer Funktion, dann wird Variablennahme nur im aktuellen C-Quelltext verwendet (wenn sich Programm aus mehreren Quelltexten zusammengesetzt wird).
\item Extern: Gegenteil von Static außerhalb einer Funktion
\item Auto: automatische Variable. Wird beim Aufruf der Funktion, die die Variablendefinition enthält, angelegt. Bei jedem Funktionsaufruf neu. Wird vernichtet, wenn Funktion beendet ist.
\end{itemize}

\section{Funktionen}
\begin{lstlisting}[language=C]
int test(){...}
// gleich wie int main: Leerer Ausgabewert ist int (nicht void!)
test(){...}
// void: unbestimmter Ausgabewert bzw. kein Ausgabewert
void test(){...}
\end{lstlisting}

\paragraph{Übung} Sinus-Funktion: $x-\frac{x^3}{3!}+\frac{x^5}{5!}-\frac{x^7}{7!}...$
\begin{lstlisting}[language=C]
#include <stdlib.h>
#include <stdio.h>
#include <math.h>

double sinus(double x);

char vbuf[128];

int main(){
	double x,y;
	fgets(vbuf,128,stdin); 
	x = atof(vbuf);
	y = sinus(x);
	printf("sin(%lf): %lf\n",x,y);
	return 0;
}

double sinus(double x){
	int i=3, vz=-1;
	double erg=x, summand=x;
	while (summand> 0.00005){
		summand = summand * x * x/(i*(i+1));
		i += 2;
		erg += summand * vz;
		vz += -1;
	}
	return erg;
}

\end{lstlisting}

\section{Pointer}

\begin{lstlisting}[language=C]
#include <stdlib.h>
#include <stdio.h>
#include <ctype.h>

char* upperstr(char* data){
	int i = 0;	
	while(data[i]!='\n'){	// (*data!='\n')
		// data[i]=toupper(data[i]);
		// Alternative zu toupper: 
		// A ist 0x41 = 0100 0001
		// a ist 0x61 = 0110 0001
		// also bloß ein Bit verschieben!
		if (data[i]>='a' && data[i]<='z'){
			data[i] &= ~(1<<5); 
			// 1101 1111, damit verunden=> alle werden negiert 0010 0000
			// &= Bitweise addition => invertierung von der 3. Stelle
		}
		// oder auch (entsprechend angepasst ohne i in while usw.)
		// *data = toupper(*data);
		// data++;
		i++;
	}
	return data;
}

char vbuf[128]

int main(){
	printf("Eingabe: ");
	fgets (vbuf,128,stdin);
	upperstr(vbuf);
	puts(vbuf);
	return 0;
}

\end{lstlisting}
Weiterführend:
\begin{lstlisting}[language=C]
#include <stdio.h>
#include <stdlib.h>

// Interpretation als Array
int mystrlen1(char *p){
	int i;
	for (i=0; p[i]!=0; i++);
	return i;
}

// Nutzen des Pointers
int mystrlen2(char *p){
	int count;
	while (*p++)
		count++;
	return count;
}

// Noch mehr nutzen des Pointers
int mystrlen3(char *p){
	char *px=p; // Pointer auf das erste Zeichen merken
	while (*p++); // Pointer hoch zählen (wenn letzte Stelle, ist *p++ 0, also false)
	return p-px-1; // (letzte Stelle)-(erste Stelle)-1
}

int intarr[] = {5,7,2,8,9};

// Interpretation als Array
int containsint(int n, int* pdata, int test){
	int i;
	for (i=0; i<n && pdata[i]!=test; i++);
	return pdata[i]==test;
}

// Nutzen der Pointer: *(pdata+i) ist das selbe wie pdata[i]
int containsInt2(int n, int* pdata, int test){
	while(*(pdata+--n)!= test && n!=0);
	return *(pdata+n)== test;
}

int mystrcmp(char* p1, char* p2){
	int i;
	for (i=0; p1[i]==p2[i] && p1[i]; i++);
	return p1[i]-p2[i];
}

int mystrcmp2(char* p1, char* p2){
	while (*p1-*p2 && *p1) p1++,p2++;
	return *p1-*p2;
}

int main (int argc, char* argv[]){
	int i, len;
	
	i=atoi(argv[1]);
	if(containsInt(sizeof intarr/sizeof(int), intarr, i))	puts("Enthalten");
	else																									puts("nicht Enthalten");
	
	if(mystrcmp(argv[1],argv[2])==0)	printf("%s gleich %s\n",argv[1],argv[2]);
	else															printf("%s ungleich %s\n",argv[1],argv[2]);
	
	for (i=0; i<argc; i++){
		puts(argv[i]);	// Ausgabe der Eingabeparameter
										// (wenn aufgerufen durch "./a.out e1 e2" werden
										// "./a.out", "e1" und "e2" ausgegeben)
		// Ausgabe des jeweils ersten Zeichens
		printf("%c\n".argv[i][0]);	// alternativ auch *argv[i]
		// argv[0] ist immer der Programmname
		
		for (i=0; i<argc; i++){
			len=mystrlen1(argv[i]);
			printf("len: %d\n", len);
		}
	}
	return 0;
}
\end{lstlisting}

\section{Benutzerdefinierte Datentypen}

Enum:

Aufzählungstyp (festgesetzte Bezeichnungen auf einen integer-Wert).\medskip\\
Struct:

Zusammenfassung von mehreren Komponenten (unterschiedliche eingebaute Dateitypen als un- intialisierte Variablen), die durch einen Namen beschrieben werden. Verwendung zur Modellierung eines Sachverhalts (wie im Beispiel Student mit seinen Eigenschaften).\medskip\\
Typedef:

Es wird ein synonymer Typname für einen existierenden Typnamen erstellt. So kann die Variable- ninitialisierung verkürzt werden (im Skript: struct tStudent$\rightarrow$tStud ).\medskip\\
Union:

Datensätze werden im Vergleich zum Struct übereinander geschrieben.

\section{Listen}
\paragraph{Implementation 1}(Ringliste)
\begin{tikzpicture}[scale=.4]
\draw  (0,4) rectangle (6,6);
\node [right] at (0,5) {first};
\draw  (0,2)rectangle (6,4);
\node [right] at (0,3) {last};
\draw  (0,0) rectangle (6,2);
\node at (0,1) [right] {current};
\draw  (11,6) rectangle (17,4);
\draw  (11,4) rectangle (17,2) node[pos=.5]{Max};
\draw  (22,6) rectangle (28,4);
\draw  (22,4) rectangle (28,2) node[pos=.5]{Moritz};

\draw [decorate, decoration={brace, amplitude=5pt}] (6,-1) -- (0,-1) node[pos=.5,below=.1]{$tList$};
\draw [decorate, decoration={brace, amplitude=5pt}] (17,-1) -- (11,-1) node[pos=.5,below=.1]{$TCNCT$};
\fill (4,5) circle[radius=.2];
\fill (25,5) circle[radius=.2];
\fill (14,5) circle[radius=.2];
\fill (4,1) circle[radius=.2];
\fill (4,3) circle[radius=.2];
\draw [-latex] (5,5) -- (13,5);
\draw [-latex] plot[smooth, tension=.7] coordinates {(5,3) (8,4) (11,7) (19.5,7)  (24,5.5)};
\draw [-latex] (15,5) -- (24,5);
\end{tikzpicture}\\
list.h:
\begin{lstlisting}[language=C]
// Strukturtyp für Konnektor (Element mit Inhalt):
typedef struct TCNCT{
	struct TCNCT* next;	// tCnct geht noch nicht innerhalb!
	void *pItem;	// void für generische Daten
}tCnct;

typedef struct{
	tCnct* pFirst;
	tCnct* pLast;
	tCnct* tCurr;
}tList;
\end{lstlisting}
Listenimplementation (list.c):
\begin{lstlisting}[language=C]
#include <stlib.h>
#include "list.h"

// erzeugt leere Liste:
tList *CreateList(void){
	tList* ptmp;
	// Speicher freigeben:
	ptmp=malloc(sizeof(tList));
	if(ptmp!=NULL){
		// offene Liste: anfängliches tList hat nur NULL-Pointer
		ptmp->pFirst=ptmp->pLast=ptmp->pCurr=NULL;
	}
	return ptmp;
}

// hinten einfügen:
int InsertTail (tList* pList, void *pItemIns){
	// Verschieden Situationen: Anfügen an leere oder schon vorhandene Liste
	tCnct *ptmp = malloc(sizeof(tCnct));
	ptmp->next=NULL;
	if(ptmp){
		ptmp->pItem = pItemIns;	// Connector mit Inhalt füllen
		if (pList->pFirst!=NULL){	// Liste Leer
			pList->pFirst=pList->pLast = ptmp;
		} else {	// Liste enthält schon Konnektoren
			pList->pLast->next=ptmp;	// Das vorher letzte Element zeigt nun auf das eingefügte
																// und damit neue letzte Element
			pList->pLast = ptmp;			// das neue letzte Element
		}
		pList->pCurr=ptmp;	// Das Element, mit dem zuletzt hantiert wurde ist pCurr
	}
	return (int)ptmp;
}

// gibt ersten Eintrag aus:
void* GetFirst (tList* pList){
	tCnct *ptmp;
	ptmp = pList->pFirst;
	if (ptmp){
		pList->pCurr=ptmp;
		return ptmp->pItem;
	}
	return NULL;
}

// gibt nächsten Eintrag aus:
void* GetNext (tList* pList){
	tCnct *ptmp = pList->pCurr;
	if (ptmp){
		if(ptmpt==pList->pLast){ // kein Nachfolger vorhanden
			return NULL;
		} else {
			ptmp = ptmp->next;
			pList->pCurr = ptmp;
			return ptmp->iItem;
		}
	} 
	return GetFirst(pList);
}

// gibt letzten Eintrag aus:
void* GetLast (tList* pList){
}

\end{lstlisting}

Vorgehen bei malloc/fopen usw. immer:
\begin{enumerate}
\item Malloc machen
\item Überprüfen, ob es geklappt hat!
\end{enumerate}

Unterschied: Offene Liste und Ringliste. Offene Liste startet mit NULL-Zeigern.

\paragraph{Implementation 2} (Doppelt-verkettete Ringlist)\\
\begin{tikzpicture}[scale=.4]
\draw  (0,4) rectangle (6,6);
\node [right] at (0,5) {next};
\draw  (0,2)rectangle (6,4);
\node [right] at (0,3) {previous};
\draw  (0,0) rectangle (6,2) node[pos=.5]{};
\fill (4.5,5) circle[radius=.2];
\fill (25,5) circle[radius=.2];
\fill (14,5) circle[radius=.2];
\fill (4.5,3) circle[radius=.2];
\fill (14,3) circle[radius=.2];
\fill (25,3) circle[radius=.2];
\draw [-latex] (5.5,5) -- (13,5);
\draw [-latex] (15,5) -- (24,5);
\draw  (11,6) rectangle (17,4);
\draw  (11,4) rectangle (17,2);
\draw  (11,2) rectangle (17,0) node[pos=.5]{Hans};
\draw  (22,6) rectangle (28,4);
\draw  (22,4) rectangle (28,2);
\draw  (22,2) rectangle (28,0) node[pos=.5]{Anna};
\draw [-latex] (26,5) -- (29,5) -- (29,7) -- (-2,7) -- (-2,5) -- (-1,5);
\draw [-latex] (13,3) -- (5.5,3);
\draw [-latex] (24,3) -- (15,3);
\draw [-latex] (-1,3) -- (-2,3) -- (-2,-1) -- (29,-1) -- (29,3) -- (26,3);
\end{tikzpicture}\\
Vorteil: Jedes Element hat einen Vorgänger und einen Nachfolger. Dadurch reicht eine Funktion, die nach einem Element ein neues einfügen kann. Das kann an beliebiger Stelle passieren.\\
list.h\\

\section{Binäre Bäume}

\begin{tikzpicture}[scale=.2]
\draw  (0,0) rectangle (3,-3);
\draw  (3,0) rectangle (6,-3);
\draw  (0,3) rectangle (6,0) node[pos=.5]{Franz};
\draw  (5,-8) rectangle (8,-11);
\draw  (8,-8) rectangle (11,-11);
\draw  (5,-5) rectangle (11,-8) node[pos=.5]{Paul};
\draw  (11,-16) rectangle (14,-19);
\draw  (14,-16) rectangle (17,-19);
\draw  (11,-13) rectangle (17,-16) node[pos=.5]{Max};
\draw  (1,-16) rectangle (4,-19);
\draw  (4,-16) rectangle (7,-19);
\draw  (1,-13) rectangle (7,-16) node[pos=.5]{Peter};
\draw  (-8,-8) rectangle (-5,-11);
\draw  (-5,-8) rectangle (-2,-11);
\draw  (-8,-5) rectangle (-2,-8) node[pos=.5]{Felix};
\draw  (-12,-16) rectangle (-9,-19);
\draw  (-9,-16) rectangle (-6,-19);
\draw  (-12,-13) rectangle (-6,-16) node[pos=.5]{Erik};
\draw [-latex] (1.5,-1.5) -- (-5,-4.5);
\draw [-latex] (4.5,-1.5) -- (8,-4.5);
\draw [-latex] (6.5,-9.5) -- (4,-12.5);
\draw [-latex] (9.5,-9.5) -- (14,-12.5);
\draw [-latex] (-6.5,-9.5) -- (-9,-12.5);
\end{tikzpicture}

\begin{lstlisting}[language=C]
#include <stdio.h>
#include <string.h>
#include <stdlib.h>

typedef struct{
	char *pdata;
	// pointer weil: 
	// char geht nicht: nur ein character
	// string gibt's noch nicht
	// char[] ungünstig, weil die Inhalte auch unterschiedlich lang seien könnten.
	// also char*
	struct TNODE *pl;
	struct TNODE *pr;
} tNode

tNode treeInit={};

char* data[]={"moritz", "max", "melli", "albert", "hans", "peter", "malte",
							"maximilian", "james", "johannes" , "paul",  NULL);
tNode *pTree;
char buf[128];

void addToTree(tNode *pt, char* pdata){
	if(pt->pdata == NULL){ // im Falle, dass noch kein Node existiert
		pt->pdata=pdata;
	} else {
		if (strcmp(pt->pdata, pdata)>0){
			//links
			if(pt->pl == NULL){
				pt->pl = malloc(sizeof(tNode));
				*(pt->pl) = treeInit;
				attToTree(pt->pl,pdata);
			}
		} else {
			//recht
			if(pt->pr == NULL){
				pt->pr = malloc(sizeof(tNode));
				*(pt->pr) = treeInit;
				attToTree(pt->pl,pdata);
			}
		}
	}
}

void dispTree(tNode *pt){
	if (pt->pl != NULL) dispTree(pt->pl);
	puts(pt-pdata);
	if (pt->pr != NULL) dispTree(pt->pr);
}

int main(){
	char *ptmp;
	char **p=data;
	// erstes Node erstellen (leer):
	pTree = malloc(sizeof(tNode));
	*pTree = treeInit;
	
	while  (*p){
		addToTree(pTree, *p); 
		p++;
	}
	
	dispTree(pTree);
}
\end{lstlisting}

Alternativ mit array, void pdata (Funktionspointer) usw.:
\begin{lstlisting}[language=C]
#include <stdio.h>
#include <string.h>
#include <stdlib.h>

typedef struct{
	void *pdata;
	struct TNODE* px[2];
} tNode

tNode treeInit={};

char* data[]={"moritz", "max", "melli", "albert", "hans", "peter", "malte",
							"maximilian", "james", "johannes" , "paul",  NULL);
tNode *pTree;
char buf[128];

int mycmp(void*p1, void*p2){
	return (strcmp((char*)p1, (char*)p2>0)?0:1;
}

void addToTree(tNode *pt, void* pdata,  int (*fcmp)(void*, void*)){
	int i;
	if(pt->pdata == NULL){
		pt->pdata=pdata;
	} else {
		i=fcmp(pt->pdata, pdata);
			if(pt->px[i] == NULL){
				pt->px[i] = malloc(sizeof(tNode));
				*(pt->px[i]) = treeInit;
			}
			attToTree(pt->px[i],pdata, mycmp);
	}
}

void mydisp(void *pdata){
	puts((char*)pdata);
}

void dispTree(tNode *pt, void (*putx)(void*)){
	if (pt->px[0] != NULL) dispTree(pt->px[0], putx);
	putx(pt-pdata);
	if (pt->px[1] != NULL) dispTree(pt->px[1], putx);
}

int main(){
	char *ptmp;
	char **p=data;
	// erstes Node erstellen (leer):
	pTree = malloc(sizeof(tNode));
	*pTree = treeInit;
	
	while  (*p){
		addToTree(pTree, *p, mycmp); 
		p++;
	}
	
	dispTree(pTree, mydisp);

	while (1){
		fgets(buf,128,stdin);
		buf[strlen(buf)-1]=0;
		ptmp=malloc(strlen(buf)+1);
		strcpy(ptmp,buf);
		puts("==========");
		addToTree(pTree, ptmp, mycmp);
		dispTree(pTree, mydisp);
}
\end{lstlisting}

\newpage
%\printbibliography
\end{document}